\documentclass{article}

\usepackage{graphicx}
\usepackage{booktabs}
\usepackage{multirow}
\usepackage{xcolor}
\usepackage{float}
\usepackage{placeins}
\usepackage{algorithm}
\usepackage{algpseudocode}
\usepackage{amsmath,amssymb}
\algnewcommand{\LineComment}[1]{\State \(\triangleright\) \textit{#1}}
\begin{document}


% ============================================================
% Latent-space GRPO training algorithm
% ============================================================
\begin{algorithm}[t]
\caption{PhyPO-T2V with Best-of-\(K\) Trajectory Selection}
\label{alg:phypo-t2v}
\footnotesize
\setlength{\algorithmicindent}{1em}
\begin{algorithmic}[1]

\Require Prompt \(p\), T2V model \(f_\theta\), decoder \(\mathcal D\),
reward \(R\), steps \(T\), training window
\(\mathcal T_{\mathrm{train}}\), rollouts \(K\), rollout noise
\(\sigma_{\mathrm{roll}}\), policy variance \(\sigma^2\), clip
\(\epsilon\), constant \(\epsilon_A\), learning rate \(\eta\)
\Ensure Updated parameters \(\theta\)

\State \(z_T\sim\mathcal N(0,I)\), \quad
\(\theta_{\mathrm{old}}\gets\theta\)
\State \(\mu_\phi(z,t,p)\gets
\operatorname{SolverStep}\bigl(z,f_\phi(z,t,p),t\bigr)\)

\For{\(t=T,T-1,\ldots,1\)}
    \If{\(t\notin\mathcal T_{\mathrm{train}}\)}
        \State \(z_{t-1}\gets
        \mu_{\theta_{\mathrm{old}}}(z_t,t,p)\);
        \textbf{continue}
    \EndIf

    \For{\(k=1,\ldots,K\)}
        \State \(\varepsilon^{(k)}\sim\mathcal N(0,I)\), \quad
        \(z_t^{(k)}
        \gets z_t+\sigma_{\mathrm{roll}}\varepsilon^{(k)}\)

        \State \(a_{\mathrm{old}}^{(k,t)}
        \gets f_{\theta_{\mathrm{old}}}(z_t^{(k)},t,p)\)

        \State \(\hat{x}_0^{(k,t)}
        \gets\operatorname{EstimateX0}
        (z_t^{(k)},a_{\mathrm{old}}^{(k,t)},t)\)

        \State \(\hat v_0^{(k,t)}
        \gets\mathcal D(\hat x_0^{(k,t)})\), \quad
        \(R_t^{(k)}\gets R(\hat v_0^{(k,t)},p)\)

        \State \(z_{t-1}^{(k)}
        \sim\pi_{\theta_{\mathrm{old}}}
        (\cdot\mid z_t^{(k)},p,t)
        =
        \mathcal N\!\left(
        \mu_{\theta_{\mathrm{old}}}(z_t^{(k)},t,p),
        \sigma^2I
        \right)\)
    \EndFor

    \State \(k_t^\star\gets\arg\max_k R_t^{(k)}\), \quad
    \(z_{t-1}\gets z_{t-1}^{(k_t^\star)}\)
    \Comment{advance with best rollout}

    \State \(\mu_{R,t}\gets K^{-1}\sum_{j=1}^K R_t^{(j)}\), \quad
    \(\sigma_{R,t}\gets
    \sqrt{K^{-1}\sum_{j=1}^K
    (R_t^{(j)}-\mu_{R,t})^2}\)

    \For{\(k=1,\ldots,K\)}
        \State \(A_t^{(k)}
        \gets\operatorname{sg}\!\left[
        \frac{R_t^{(k)}-\mu_{R,t}}
        {\sigma_{R,t}+\epsilon_A}
        \right]\)

        \State \(\ell_\theta^{(k,t)}
        \gets\log\pi_\theta
        (z_{t-1}^{(k)}\mid z_t^{(k)},p,t)\)

        \State \(\ell_{\mathrm{old}}^{(k,t)}
        \gets\log\pi_{\theta_{\mathrm{old}}}
        (z_{t-1}^{(k)}\mid z_t^{(k)},p,t)\)

        \State \(\rho_{k,t}(\theta)
        \gets\exp\!\left(
        \ell_\theta^{(k,t)}
        -\ell_{\mathrm{old}}^{(k,t)}
        \right)\)
    \EndFor
\EndFor

\State \(\displaystyle
\mathcal L_{\mathrm{GRPO}}\gets
-\frac{1}{K|\mathcal T_{\mathrm{train}}|}
\sum_{t\in\mathcal T_{\mathrm{train}}}
\sum_{k=1}^K
\min\!\left(
\rho_{k,t}(\theta)A_t^{(k)},
\operatorname{clip}
(\rho_{k,t}(\theta),1-\epsilon,1+\epsilon)
A_t^{(k)}
\right)\)

\State \(\theta\gets
\theta-\eta\nabla_\theta\mathcal L_{\mathrm{GRPO}}\)

\end{algorithmic}
\end{algorithm}
\FloatBarrier
\paragraph{Intermediate $x_0$ reward prediction.}
At each intermediate denoising timestep $t$, we use the model's
$x_0$-prediction $\hat{x}_0^{(k,t)}$ as an estimate of the terminal clean
video latent. We decode $\hat{x}_0^{(k,t)}$ into the complete predicted
temporal video $\hat{v}_0^{(k,t)} = D(\hat{x}_0^{(k,t)})$ and compute the
reward as $R_t^{(k)} = R(\hat{v}_0^{(k,t)}, p)$. Thus, the reward evaluates
the predicted terminal video trajectory rather than the local latent
transition $z_t^{(k)} \rightarrow z_{t-1}^{(k)}$ alone.

\paragraph{Connection to the GRPO objective.}
Algorithm~\ref{alg:phypo-t2v} instantiates Eq.~(1) in the latent
denoising space, where the rollout index $i$ corresponds to $k$ and the optimization
step $j$ corresponds to the denoising time step $t$. We define the policy action
as the next latent $z_{t-1}^{(k)}$ and the policy state as
$s_t^{(k)}=(z_t^{(k)},p,t,s_t^{\mathrm{solver}})$. Consequently, the
importance ratio in Eq.~(1) becomes
\[
\rho_{k,t}(\theta)
=
\frac{
\pi_\theta(z_{t-1}^{(k)} \mid z_t^{(k)},p,t)
}{
\pi_{\theta_{\mathrm{old}}}
(z_{t-1}^{(k)} \mid z_t^{(k)},p,t)
}.
\]
We model the latent transition policy as
$\pi_\theta(z_{t-1}\mid z_t,p,t)
=\mathcal{N}(z_{t-1};\mu_\theta(z_t,p,t),\sigma_t^2 I)$,
where $\mu_\theta$ is the next-latent prediction produced by the adapter.
Thus,
\[
\log \pi_\theta(z_{t-1}^{(k)} \mid z_t^{(k)},p,t)
=
-\frac{
\|z_{t-1}^{(k)}-\mu_\theta(z_t^{(k)},p,t)\|_2^2
}{
2\sigma_t^2
}
+C_t.
\]

The reward is not evaluated on the local transition
$z_t^{(k)}\rightarrow z_{t-1}^{(k)}$ alone. At each intermediate
denoising timestep $t$, the model predicts a terminal clean-video latent
$\hat{x}_0^{(k,t)}$, which is decoded into the complete predicted temporal
video $\hat{v}_0^{(k,t)}=D(\hat{x}_0^{(k,t)})$. We compute
$R_t^{(k)}=R(\hat{v}_0^{(k,t)},p)$ and normalize the $K$ rewards obtained
at the same timestep to produce the group-relative advantage
$A_t^{(k)}$. Consequently, Eq.~(1) is instantiated using
$\rho_{k,t}(\theta)$ and $A_t^{(k)}$: policy optimization is local to the
intermediate denoising state, whereas reward supervision evaluates the
predicted terminal video trajectory.

\paragraph{Latent-Space GRPO.}
We instantiate GRPO over the latent denoising process by treating each
intermediate denoising state as
$s_t^{(k)}=(z_t^{(k)},p,t,s_t^{\mathrm{solver}})$ and the corresponding
policy action as the next latent $a_t^{(k)}=z_{t-1}^{(k)}$.
Accordingly, the generic GRPO importance ratio in Eq.~(1) becomes
\begin{equation}
\rho_{k,t}(\theta)
=
\frac{
\pi_\theta
\left(
z_{t-1}^{(k)}
\mid
z_t^{(k)},p,t
\right)
}{
\pi_{\theta_{\mathrm{old}}}
\left(
z_{t-1}^{(k)}
\mid
z_t^{(k)},p,t
\right)
}.
\end{equation}
We model the latent transition policy as a Gaussian,
\begin{equation}
\pi_\theta(z_{t-1}\mid z_t,p,t)
=
\mathcal{N}
\left(
z_{t-1};
\mu_\theta(z_t,p,t),
\sigma_t^2 I
\right),
\end{equation}
where $\mu_\theta(z_t,p,t)$ denotes the next-latent prediction of the
video generation model. The corresponding policy log-likelihood is
\begin{equation}
\log \pi_\theta
\left(
z_{t-1}^{(k)}
\mid
z_t^{(k)},p,t
\right)
=
-\frac{
\left\|
z_{t-1}^{(k)}
-
\mu_\theta(z_t^{(k)},p,t)
\right\|_2^2
}{
2\sigma_t^2
}
+C_t.
\end{equation}

\paragraph{Intermediate $x_0$ Reward Prediction.}
Rather than evaluating the local latent transition
$z_t^{(k)}\rightarrow z_{t-1}^{(k)}$ directly, we use the model's
$x_0$-prediction at each intermediate denoising timestep as an estimate
of the terminal clean video latent,
\begin{equation}
\hat{x}_0^{(k,t)}
=
f_{\theta_{\mathrm{old}}}
\left(
z_t^{(k)},t,p
\right).
\end{equation}
We decode this prediction into the complete predicted temporal video
trajectory,
\begin{equation}
\hat{v}_0^{(k,t)}
=
D\left(\hat{x}_0^{(k,t)}\right),
\end{equation}
and compute
\begin{equation}
R_t^{(k)}
=
R\left(\hat{v}_0^{(k,t)},p\right).
\end{equation}
Thus, reward supervision evaluates the predicted terminal video
trajectory, while policy optimization remains local to the intermediate
denoising state.

For the $K$ rollouts evaluated at timestep $t$, we compute the
group-relative advantage as
\begin{equation}
A_t^{(k)}
=
\frac{
R_t^{(k)}-\mu_{R,t}
}{
\sigma_{R,t}+\epsilon_A
},
\end{equation}
where
\begin{equation}
\mu_{R,t}
=
\frac{1}{K}\sum_{k=1}^{K}R_t^{(k)},
\qquad
\sigma_{R,t}
=
\sqrt{
\frac{1}{K}
\sum_{k=1}^{K}
\left(R_t^{(k)}-\mu_{R,t}\right)^2
}.
\end{equation}
The resulting $A_t^{(k)}$ and $\rho_{k,t}(\theta)$ are substituted
directly into the clipped GRPO objective in Eq.~(1). Overall, the
optimization follows
\begin{equation}
z_t^{(k)}
\rightarrow
\hat{x}_0^{(k,t)}
\rightarrow
D(\hat{x}_0^{(k,t)})
\rightarrow
R_t^{(k)}
\rightarrow
A_t^{(k)}
\rightarrow
\rho_{k,t}(\theta)
\rightarrow
\mathcal{L}_{\mathrm{GRPO}},
\end{equation}
providing full-video physics feedback at intermediate denoising
timesteps without requiring a complete denoising rollout from each
intermediate state.



\begin{table}[H]
\centering
\small
\setlength{\tabcolsep}{6pt}
\renewcommand{\arraystretch}{1.15}
\begin{tabular}{lcccc}
\toprule
\textbf{Model}
& \textbf{Option} $\uparrow$
& \textbf{Formula} $\uparrow$
& \textbf{Substitution} $\uparrow$
& \textbf{F-T} $\uparrow$ \\
\midrule
Wan2.1-1.3B       & -- & -- & -- & -- \\
\textbf{+ PhyPO}  & \textbf{--} & \textbf{--} & \textbf{--} & \textbf{--} \\
\midrule
LTX-Video-2B      & -- & -- & -- & -- \\
\textbf{+ PhyPO}  & \textbf{--} & \textbf{--} & \textbf{--} & \textbf{--} \\
\midrule
Cosmos-2B         & -- & -- & -- & -- \\
\textbf{+ PhyPO}  & \textbf{--} & \textbf{--} & \textbf{--} & \textbf{--} \\
\bottomrule
\end{tabular}
\caption{\textbf{Physics formulation results on the adapted
Apple-PI F-T benchmark.} Initial physical quantities originally
provided through image annotations are instead included directly
in the text prompt for T2V evaluation.}
\label{tab:apple-pi-ft}
\end{table}

% ============================================================
% Human evaluation: physical reasoning
% ============================================================

% Table 1
\begin{table}[H]
\centering
\small
\renewcommand{\arraystretch}{1.7}
\setlength{\tabcolsep}{3pt}
\resizebox{\textwidth}{!}{%
\begin{tabular}{lcccccccccc}
\toprule
\textbf{Video Backbone} &
\textbf{Falling} & \textbf{Projectile} & \textbf{Swinging} & \textbf{Spinning} &
\textbf{Compressing} & \textbf{Sliding} & \textbf{Bouncing} & \textbf{Flowing} &
\textbf{Floating} & \textbf{Splashing} \\
\midrule
LTX-Video-2B &
& & & & & & & & & \\
LTX-Video-13B &
& & & & & & & & & \\
Cosmos-Predict2.5-2B &
& & & & & & & & & \\
Cosmos-Predict2.5-14B &
& & & & & & & & & \\
Wan2.1-1.3B &
& & & & & & & & & \\
Wan2.1-14B &
& & & & & & & & & \\
\bottomrule
\end{tabular}%
}
\caption{Human evaluation preference win rate (\%) of physical reasoning across video generation models.}
\label{tab:human-physics-eval}
\end{table}
\FloatBarrier


% ============================================================
% Cross-VLM evaluation
% ============================================================

% Table 2
\begin{table}[H]
\centering
\small
\renewcommand{\arraystretch}{1.5}
\setlength{\tabcolsep}{4pt}
\resizebox{\textwidth}{!}{%
\begin{tabular}{llcccccccccc}
\toprule
\textbf{Video Backbone} & \textbf{Evaluator} &
\textbf{Falling} & \textbf{Projectile} & \textbf{Swinging} & \textbf{Spinning} &
\textbf{Compressing} & \textbf{Sliding} & \textbf{Bouncing} & \textbf{Flowing} &
\textbf{Floating} & \textbf{Splashing} \\
\midrule

\multirow{2}{*}{LTX-Video-2B}
& Qwen2-VL-72B-Instruct
& \textbf{\textcolor{blue}{59.1}}
& \textbf{\textcolor{blue}{63.1}}
& \textbf{\textcolor{blue}{65.1}}
& \textbf{\textcolor{blue}{66.3}}
& \textbf{\textcolor{blue}{52.0}}
& \textbf{\textcolor{blue}{54.7}}
& \textbf{\textcolor{blue}{67.6}}
& \textbf{\textcolor{blue}{69.5}}
& \textbf{\textcolor{blue}{72.9}}
& \textbf{\textcolor{blue}{70.2}} \\
& Gemini-3.1-Flash-Lite
& \textbf{\textcolor{red}{46.0}}
& \textbf{\textcolor{blue}{55.4}}
& \textbf{\textcolor{blue}{58.0}}
& \textbf{\textcolor{blue}{62.0}}
& \textbf{\textcolor{red}{47.1}}
& \textbf{\textcolor{red}{49.7}}
& \textbf{\textcolor{blue}{52.8}}
& \textbf{\textcolor{blue}{58.5}}
& \textbf{\textcolor{blue}{61.1}}
& \textbf{\textcolor{blue}{54.3}} \\
\midrule

\multirow{2}{*}{LTX-Video-13B}
& Qwen2-VL-72B-Instruct
& \textbf{\textcolor{blue}{70.0}}
& \textbf{\textcolor{blue}{71.4}}
& \textbf{\textcolor{blue}{70.3}}
& \textbf{\textcolor{blue}{77.4}}
& \textbf{\textcolor{blue}{80.9}}
& \textbf{\textcolor{blue}{64.9}}
& \textbf{\textcolor{blue}{76.4}}
& \textbf{\textcolor{blue}{86.3}}
& \textbf{\textcolor{blue}{73.7}}
& \textbf{\textcolor{blue}{82.0}} \\
& Gemini-3.1-Flash-Lite
& \textbf{\textcolor{blue}{51.4}}
& \textbf{\textcolor{red}{48.0}}
& \textbf{\textcolor{red}{44.6}}
& \textbf{\textcolor{red}{44.0}}
& \textbf{\textcolor{blue}{63.4}}
& \textbf{\textcolor{red}{44.0}}
& \textbf{\textcolor{red}{43.7}}
& \textbf{\textcolor{red}{43.1}}
& \textbf{\textcolor{red}{46.6}}
& \textbf{\textcolor{blue}{52.2}} \\
\midrule

\multirow{2}{*}{Cosmos-Predict2.5-2B}
& Qwen2-VL-72B-Instruct
& \textbf{\textcolor{blue}{67.4}}
& \textbf{\textcolor{blue}{77.4}}
& \textbf{\textcolor{blue}{70.6}}
& \textbf{\textcolor{blue}{67.1}}
& \textbf{\textcolor{blue}{64.6}}
& \textbf{\textcolor{blue}{56.0}}
& \textbf{\textcolor{blue}{64.1}}
& \textbf{\textcolor{blue}{59.9}}
& \textbf{\textcolor{blue}{51.4}}
& \textbf{\textcolor{blue}{71.4}} \\
& Gemini-3.1-Flash-Lite
& \textbf{\textcolor{blue}{60.9}}
& \textbf{\textcolor{blue}{72.0}}
& \textbf{\textcolor{blue}{64.3}}
& \textbf{\textcolor{blue}{62.9}}
& \textbf{\textcolor{blue}{58.9}}
& \textbf{\textcolor{blue}{58.6}}
& \textbf{\textcolor{blue}{67.1}}
& \textbf{\textcolor{blue}{61.1}}
& \textbf{\textcolor{blue}{56.0}}
& \textbf{\textcolor{blue}{64.9}} \\
\midrule

\multirow{2}{*}{Cosmos-Predict2.5-14B}
& Qwen2-VL-72B-Instruct
& & & & & & & & & & \\
& Gemini-3.1-Flash-Lite
& & & & & & & & & & \\
\midrule

\multirow{2}{*}{Wan2.1-1.3B}
& Qwen2-VL-72B-Instruct
& \textbf{\textcolor{blue}{56.3}}
& \textbf{\textcolor{blue}{50.3}}
& \textbf{\textcolor{blue}{54.0}}
& \textbf{\textcolor{red}{46.6}}
& \textbf{\textcolor{blue}{63.7}}
& \textbf{\textcolor{red}{41.4}}
& \textbf{\textcolor{red}{48.1}}
& \textbf{\textcolor{red}{47.1}}
& \textbf{\textcolor{blue}{51.7}}
& \textbf{\textcolor{blue}{50.3}} \\
& Gemini-3.1-Flash-Lite
& \textbf{\textcolor{blue}{53.1}}
& \textbf{\textcolor{blue}{55.7}}
& \textbf{\textcolor{red}{45.1}}
& \textbf{\textcolor{red}{48.9}}
& \textbf{\textcolor{blue}{55.7}}
& \textbf{\textcolor{red}{40.6}}
& \textbf{\textcolor{red}{48.4}}
& \textbf{\textcolor{red}{45.7}}
& \textbf{\textcolor{red}{46.9}}
& \textbf{\textcolor{red}{47.7}} \\
\midrule

\multirow{2}{*}{Wan2.1-14B}
& Qwen2-VL-72B-Instruct
& & & & & & & & & & \\
& Gemini-3.1-Flash-Lite
& & & & & & & & & & \\
\bottomrule
\end{tabular}%
}
\caption{Cross-VLM preference win rate (\%) of physics-awareness post-trained models over baselines across six video generation backbones. Values above 50\% are shown in bold blue; values at or below 50\% are shown in bold red.}
\label{tab:human_vlm_physics_eval}
\end{table}
\FloatBarrier


% ============================================================
% PAI-Bench: Physics category
% ============================================================

% Table 3
\begin{table}[H]
\centering
\small
\renewcommand{\arraystretch}{1.3}
\setlength{\tabcolsep}{8pt}
\begin{tabular}{lc}
\toprule
\textbf{Model} & \textbf{Physics} $\uparrow$ \\
\midrule
Wan2.1-T2V-1.3B
& 84.8 \\
+ PhyPO
& \textbf{\textcolor{blue}{85.1}} \\
\midrule
LTX-Video-2B
& 84.8 \\
+ PhyPO
& \textbf{\textcolor{blue}{87.7}} \\
\midrule
Cosmos-predict2.5-2B
& 91.7 \\
+ PhyPO
& -- \\
\bottomrule
\end{tabular}
\caption{\textbf{PAI-Bench Physics-category results.}
Higher is better. We report the Physics (PH) category because PhyPO specifically targets physical consistency.}
\label{tab:pai_bench_physics}
\end{table}
\FloatBarrier


% ============================================================
% VBench (Table 4: forced immediately after Table 3, before Table 5)
% ============================================================

% Table 4
\begin{table}[H]
\footnotesize
\setlength{\tabcolsep}{3pt}
\centering
\begin{tabular}{lcccccc}
\toprule
\textbf{Model} & SC & BC & MS & AQ & IQ & OC \\
\midrule
LTX-Video-2B
& 93.9 & 95.5 & 99.4 & 46.1 & 55.4 & 17.2 \\
{+ PhyPO}
& \textbf{\textcolor{blue}{98.6}}
& \textbf{\textcolor{blue}{98.7}}
& \textbf{\textcolor{blue}{99.6}}
& \textbf{\textcolor{blue}{52.4}}
& \textbf{\textcolor{blue}{70.2}}
& \textbf{\textcolor{blue}{21.8}} \\
\midrule
LTX-Video-13B
& \textbf{\textcolor{blue}{99.6}}
& 99.0
& \textbf{\textcolor{blue}{99.7}}
& \textbf{\textcolor{blue}{53.2}}
& 70.3
& 18.9 \\
{+ PhyPO}
& 99.2
& 99.0
& 99.6
& 52.9
& \textbf{\textcolor{blue}{73.7}}
& \textbf{\textcolor{blue}{23.1}} \\
\midrule
Cosmos-Predict2.5-2B
& 96.7 & 96.5 & 99.3 & 50.1 & \textbf{\textcolor{blue}{74.5}} & 22.6 \\
{+ PhyPO}
& \textbf{\textcolor{blue}{97.0}}
& \textbf{\textcolor{blue}{97.1}}
& 99.1
& \textbf{\textcolor{blue}{53.0}}
& 74.0
& \textbf{\textcolor{blue}{27.2}} \\
\midrule
Cosmos-Predict2.5-14B
& -- & -- & -- & -- & -- & -- \\
{+ PhyPO}
& -- & -- & -- & -- & -- & -- \\
\midrule
Wan2.1-1.3B
& 94.8
& 95.9
& 98.8
& 51.1
& 68.0
& 24.5 \\
{+ PhyPO}
& \textbf{\textcolor{blue}{95.0}}
& \textbf{\textcolor{blue}{96.6}}
& 98.8
& \textbf{\textcolor{blue}{55.2}}
& \textbf{\textcolor{blue}{70.8}}
& \textbf{\textcolor{blue}{25.8}} \\
\midrule
Wan2.1-14B
& -- & -- & -- & -- & -- & -- \\
{+ PhyPO}
& -- & -- & -- & -- & -- & -- \\
\bottomrule
\end{tabular}
\caption{\textbf{VBench results.}
SC: Subject Consistency; BC: Background Consistency; MS: Motion Smoothness;
AQ: Aesthetic Quality; IQ: Imaging Quality; OC: Overall Consistency.}
\label{tab:suppl_exp_vbench_quality}
\end{table}
\FloatBarrier
% Forced order: Table 3 (PAI) -> Table 4 (VBench) -> Table 5 (LIBERO)


% ============================================================
% Table 5: LIBERO suite
% Table 6: Spatial per-task
% then paragraph
% then Table 7: RoboTwin
% ============================================================

% --- Table 5 ---
\begin{table}[H]
\centering
\caption{LIBERO success rates (16k steps) (\%).}
\label{tab:libero-phypo-success}
\resizebox{\linewidth}{!}{%
\begin{tabular}{lcccccc}
\toprule
Task & Wan2.1-1.3B & Wan2.1-1.3B + PhyPO & $\Delta$ &
Wan2.1-14B & Wan2.1-14B + PhyPO & $\Delta$ \\
\midrule
Spatial
& 91.00 & 83.60 & \textcolor{red}{$-7.40$}
& -- & -- & -- \\

Object
& 94.80 & 98.00 & \textcolor{blue}{$+3.20$}
& -- & -- & -- \\

Goal
& 64.60 & 65.00 & \textcolor{blue}{$+0.40$}
& -- & -- & -- \\

Libero-10
& 15.20 & 24.60 & \textcolor{blue}{$+9.40$}
& -- & -- & -- \\

\midrule
Overall
& 66.40 & 67.80 & \textcolor{blue}{$+1.40$}
& -- & -- & -- \\
\bottomrule
\end{tabular}%
}
\end{table}




% --- Table 6 ---
\begin{table}[H]
\centering
\small
\caption{Per-task LIBERO-Spatial success rates (\%). Cabinet/drawer tasks drive the suite drop.}
\label{tab:libero-spatial-per-task}

\renewcommand{\arraystretch}{1.4}
\resizebox{0.95\linewidth}{!}{%
\begin{tabular}{llccc}
\toprule
ID & Task (short) & Base & +PhyPO & $\Delta$ \\
\midrule
0 & between plate and ramekin & 90 & 98 & $+8$ \\
1 & next to ramekin & 98 & 96 & $-2$ \\
2 & table center & 92 & 100 & $+8$ \\
3 & on cookie box & 100 & 100 & $0$ \\
4 & in top drawer of cabinet & 76 & 40 & \textcolor{red}{$-36$} \\
5 & on ramekin & 78 & 78 & $0$ \\
6 & next to cookie box & 100 & 100 & $0$ \\
7 & on stove & 92 & 86 & $-6$ \\
8 & next to plate & 94 & 92 & $-2$ \\
9 & on wooden cabinet & 90 & 46 & \textcolor{red}{$-44$} \\
\midrule
& \textbf{Suite mean} & \textbf{91.0} & \textbf{83.6}
& \textcolor{red}{\textbf{$-7.4$}} \\
\bottomrule
\end{tabular}%
}
\end{table}



% --- Paragraph (after Tables 5 and 6) ---
\paragraph{Downstream manipulation transfer.}
Table~\ref{tab:libero-phypo-success} reports LIBERO success rates for
Wan2.1-1.3B fine-tuned with and without PhyPO for 16k steps.
PhyPO improves overall success from $66.40\%$ to $67.80\%$ ($+1.40$).
Gains are concentrated on Object ($+3.20$; $94.80{\rightarrow}98.00$) and the
long-horizon Libero-10 suite ($+9.40$; $15.20{\rightarrow}24.60$), with a small
Goal improvement ($+0.40$).
In contrast, Spatial decreases from $91.00\%$ to $83.60\%$ ($-7.40$).
Because Spatial is already near ceiling for the baseline, this regression does
not overturn the positive overall trend driven by Object and Libero-10.

A per-task breakdown (Table~\ref{tab:libero-spatial-per-task}) shows that the
Spatial drop is localized rather than suite-wide.
Eight of ten Spatial tasks remain comparable or improve under PhyPO, especially
table-top relational pick-and-place (e.g., bowl between/next to objects, bowl
from table center).
The suite mean is dominated by two cabinet-related failures: retrieving the
bowl from the top drawer ($76\%{\rightarrow}40\%$, $-36$) and picking the bowl
from the wooden cabinet surface ($90\%{\rightarrow}46\%$, $-44$).
These tasks require precise interaction with elevated or articulated furniture,
whereas most remaining Spatial tasks are planar table-top placements where
PhyPO stays strong.

We attribute this pattern to a skill trade-off induced by the physics-oriented
video prior.
PhyPO encourages representations that better capture object dynamics and
multi-step physical interaction, which aligns with gains on Object
identity/manipulation and long-horizon Libero-10.
The same bias is less helpful---and can be harmful---for fine-grained spatial
localization on cabinets and drawers, where success hinges on geometric
precision and articulated-object affordance more than on physical dynamics.
Importantly, the Spatial regression is therefore best interpreted as a
concentrated failure mode on furniture-centric subtasks, not as a general
collapse of spatial reasoning.



% --- Table 7 ---
\begin{table}[H]
\centering
\caption{RoboTwin 2.0 success rates (\%). Each task is evaluated over 100 episodes.}
\label{tab:robotwin-phypo-success}

\renewcommand{\arraystretch}{1.4}
\resizebox{0.95\linewidth}{!}{%
\begin{tabular}{lcccccc}
\toprule
Task Group
& Wan2.1-1.3B
& Wan2.1-1.3B + PhyPO
& $\Delta$
& Wan2.1-14B
& Wan2.1-14B + PhyPO
& $\Delta$ \\
\midrule
Basic Manipulation
& -- & -- & --
& -- & -- & -- \\

Spatial Reasoning
& -- & -- & --
& -- & -- & -- \\

Object Interaction
& -- & -- & --
& -- & -- & -- \\

Long-Horizon Tasks
& -- & -- & --
& -- & -- & -- \\

\midrule
\textbf{Overall}
& \textbf{--} & \textbf{--} & \textbf{--}
& \textbf{--} & \textbf{--} & \textbf{--} \\
\bottomrule
\end{tabular}%
}
\end{table}


\begin{algorithm}[t]
\caption{PhyPO-I2V with Best-of-\(K\) Trajectory Selection}
\label{alg:phypo-i2v}
\footnotesize
\setlength{\algorithmicindent}{1em}
\begin{algorithmic}[1]

\Require Prompt \(p\), image \(I_{\mathrm{cond}}\), I2V model \(f_\theta\),
decoder \(\mathcal D\), reward \(R\), steps \(T\), training window
\(\mathcal T_{\mathrm{train}}\), rollouts \(K\), noise
\(\sigma_{\mathrm{roll}}\), policy variance \(\sigma^2\), clip
\(\epsilon\), constant \(\epsilon_A\), learning rate \(\eta\)
\Ensure Updated parameters \(\theta\)

\State \(c\gets\bigl(p,E_{\mathrm{cond}}(I_{\mathrm{cond}})\bigr)\),
\(z_T\sim\mathcal N(0,I)\), \(\theta_{\mathrm{old}}\gets\theta\)
\State \(\mu_\phi(z,t,c)\gets
\operatorname{SolverStep}\bigl(z,f_\phi(z,t,c),t\bigr)\)

\For{\(t=T,T-1,\ldots,1\)}
    \If{\(t\notin\mathcal T_{\mathrm{train}}\)}
        \State \(z_{t-1}\gets\mu_{\theta_{\mathrm{old}}}(z_t,t,c)\);
        \textbf{continue}
    \EndIf

    \For{\(k=1,\ldots,K\)}
        \State \(\varepsilon^{(k)}\sim\mathcal N(0,I)\), \quad
        \(z_t^{(k)}\gets z_t+\sigma_{\mathrm{roll}}\varepsilon^{(k)}\)
        \State \(a_{\mathrm{old}}^{(k,t)}
        \gets f_{\theta_{\mathrm{old}}}(z_t^{(k)},t,c)\)
        \State \(\hat{x}_0^{(k,t)}
        \gets\operatorname{EstimateX0}
        (z_t^{(k)},a_{\mathrm{old}}^{(k,t)},t)\)
        \State \(\hat v_0^{(k,t)}\gets\mathcal D(\hat x_0^{(k,t)})\), \quad
        \(R_t^{(k)}\gets R(\hat v_0^{(k,t)},p,I_{\mathrm{cond}})\)
        \State \(z_{t-1}^{(k)}
        \sim\pi_{\theta_{\mathrm{old}}}(\cdot\mid z_t^{(k)},c,t)
        =\mathcal N(\mu_{\theta_{\mathrm{old}}}(z_t^{(k)},t,c),\sigma^2I)\)
    \EndFor

    \State \(k_t^\star\gets\arg\max_k R_t^{(k)}\), \quad
    \(z_{t-1}\gets z_{t-1}^{(k_t^\star)}\)
    \Comment{advance with best rollout}

    \State \(\mu_{R,t}\gets K^{-1}\sum_{j=1}^K R_t^{(j)}\), \quad
    \(\sigma_{R,t}\gets
    \sqrt{K^{-1}\sum_{j=1}^K(R_t^{(j)}-\mu_{R,t})^2}\)

    \For{\(k=1,\ldots,K\)}
        \State \(A_t^{(k)}\gets\operatorname{sg}\!\left[
        (R_t^{(k)}-\mu_{R,t})/(\sigma_{R,t}+\epsilon_A)\right]\)
        \State \(\ell_\theta^{(k,t)}
        \gets\log\pi_\theta(z_{t-1}^{(k)}\mid z_t^{(k)},c,t)\)
        \State \(\ell_{\mathrm{old}}^{(k,t)}
        \gets\log\pi_{\theta_{\mathrm{old}}}
        (z_{t-1}^{(k)}\mid z_t^{(k)},c,t)\)
        \State \(\rho_{k,t}(\theta)
        \gets\exp(\ell_\theta^{(k,t)}-\ell_{\mathrm{old}}^{(k,t)})\)
    \EndFor
\EndFor

\State \(\displaystyle
\mathcal L_{\mathrm{GRPO}}\gets
-\frac{1}{K|\mathcal T_{\mathrm{train}}|}
\sum_{t\in\mathcal T_{\mathrm{train}}}\sum_{k=1}^K
\min\!\left(
\rho_{k,t}(\theta)A_t^{(k)},
\operatorname{clip}(\rho_{k,t}(\theta),1-\epsilon,1+\epsilon)A_t^{(k)}
\right)\)

\State \(\theta\gets\theta-\eta\nabla_\theta\mathcal L_{\mathrm{GRPO}}\)

\end{algorithmic}
\end{algorithm}
\FloatBarrier

BlockStackingSpecifiedOrderTask evaluates physics-aware interaction because the robot must grasp, lift, orient, and precisely place multiple blocks while accounting for gravity, contact forces, friction, collisions, and object stability. Success requires more than reaching target coordinates: each block must remain balanced in the specified order without slipping, toppling, or being displaced by later placements. The task therefore tests whether the policy can produce controlled motions and adapt to the physical consequences of sequential object manipulation.


\begin{table}[t]
\centering
\caption{PhyPO win rates (\%)}
\label{tab:phypo_win_rates}
\begin{tabular}{lccccc}
\toprule
\textbf{Judge} &
\textbf{Overall} &
\textbf{Cosmos-14B} &
\textbf{Cosmos-2B} &
\textbf{Wan-1.3B} &
\textbf{Wan-14B} \\
\midrule
Qwen2.5-VL-72B-Instruct & 60.50 & 52.71 & 74.44 & 65.73 & 49.30 \\
Qwen3-VL-32B-Instruct & 60.45 & 50.70 & 78.50 & 63.33 & 49.50 \\
\bottomrule
\end{tabular}
\end{table}


\begin{table}[t]
\centering
\caption{PhyPO win rate by physics category (\%), with ties counted as PhyPO wins.}
\label{tab:phypo_category_win_rates}
\begin{tabular}{lcc}
\toprule
\textbf{Physics Category} &
\textbf{Qwen2.5-VL-72B} &
\textbf{Qwen3-VL-32B-Instruct} \\
\midrule
Falling       & 61.00 & 61.50 \\
Projectile    & 52.50 & 61.00 \\
Swinging      & 57.00 & 57.50 \\
Compression   & 70.41 & 67.86 \\
Spinning      & 64.50 & 59.50 \\
Sliding       & 66.50 & 68.50 \\
Bouncing      & 64.80 & 62.24 \\
Fluid Flowing & 51.96 & 56.86 \\
Floating      & 60.00 & 55.00 \\
Splashing     & 56.70 & 54.64 \\
\bottomrule
\end{tabular}
\end{table}

\begin{figure}[t]
\centering
\caption{\textbf{PhyPO win rate by Basic-Physics-10 physical categories.}
Each bar is one VLM judge. The number on each bar is the PhyPO win rate (\%).
The dashed line marks chance ($50\%$).
Smaller backbones (2B / 1.3B) are on the left; larger backbones (14B / 13B) are on the right.
Mean PhyPO rate is the unweighted average of the 60 category$\times$judge bars in that panel.
Bar-wise PhyPO win-over rate is the share of those 60 bars strictly above $50\%$.}
\label{fig:phypo-all-models}
\end{figure}

\begin{table}[t]
\centering
\small
\setlength{\tabcolsep}{4pt}
\caption{\textbf{Summary of pairwise PhyPO win rates} over 6 VLM judges and 10 physics categories (60 bars per video model).
The mean is the unweighted average of the category-level rates in Fig.~\ref{fig:phypo-all-models}.
The bar-wise PhyPO win-over rate is the share of those 60 bars strictly above $50\%$.}
\label{tab:phypo-vlm-summary}
\begin{tabular}{lcc}
\toprule
Video model & Mean PhyPO rate (\%) & Win-over rate (\%) \\
\midrule
NVIDIA Cosmos-Predict2.5-2B  & 70.38 & 93.33 \\
LTX-Video 2B                 & 62.14 & 80.00 \\
Wan2.1-T2V-1.3B              & 55.04 & 70.00 \\
NVIDIA Cosmos-Predict2.5-14B & 50.74 & 50.00 \\
Wan2.1-T2V-14B               & 48.61 & 46.67 \\
LTX-Video 13B                & 48.42 & 46.67 \\
\bottomrule
\end{tabular}
\end{table}

\begin{table}[t]
\centering
\footnotesize
\setlength{\tabcolsep}{4pt}
\caption{\textbf{Per-judge overall PhyPO win rate (\%)} on Basic-Physics pairwise A/B (all categories pooled).
Forced-choice A/B; position of baseline vs.\ PhyPO is shuffled per pair.}
\label{tab:phypo-vlm-per-judge}
\begin{tabular}{lcccccc}
\toprule
Video model
  & GPT-5.6-luna
  & Qwen2.5-VL-72B
  & Qwen3-VL-32B
  & InternVL3.5-30B-A3B
  & InternVL3.5-8B
  & Molmo-72B-0924 \\
\midrule
Cosmos-Predict2.5-2B  & 79.51 & 73.02 & 76.06 & 65.72 & 72.01 & 55.17 \\
LTX-Video 2B          & 55.35 & 70.91 & 64.04 & 63.43 & 63.84 & 55.35 \\
Wan2.1-T2V-1.3B       & 51.90 & 57.11 & 54.91 & 55.11 & 61.92 & 49.30 \\
Cosmos-Predict2.5-14B & 50.90 & 48.90 & 47.09 & 48.30 & 56.31 & 52.91 \\
Wan2.1-T2V-14B        & 50.50 & 45.89 & 45.89 & 47.59 & 49.90 & 52.10 \\
LTX-Video 13B         & 48.08 & 49.70 & 45.86 & 47.62 & 50.10 & 49.09 \\
\bottomrule
\end{tabular}
\end{table}

\begin{table}[t]
\centering
\small
\setlength{\tabcolsep}{6pt}
\caption{\textbf{LIBERO success rates} for the stock policy versus PhyPO.
Each suite is evaluated over 200 rollouts. Overall is the micro-average over all 800 episodes.}
\label{tab:libero-success}
\begin{tabular}{lcc}
\toprule
Suite & Stock & PhyPO \\
\midrule
libero\_spatial     & 98.0\% (196/200) & 97.0\% (194/200) \\
libero\_object      & 98.0\% (196/200) & \textbf{99.0\%} (198/200) \\
libero\_goal        & 95.0\% (190/200) & \textbf{96.0\%} (192/200) \\
libero\_10 (long)   & 91.5\% (183/200) & 91.5\% (183/200) \\
\midrule
Overall             & 95.62\% (765/800) & \textbf{95.88\%} (767/800) \\
\bottomrule
\end{tabular}
\end{table}
\end{document}
